\section*{Chapter \ref{ch:ml:large-language-models-in-finance} --- Large Language Models in Finance}

\begin{solution}{exo:ml:large-language-models-in-finance:1}
Only July to December of the second year, after the cut-off, can measure judgement, and only if the provider's cut-off is
right and the model was not updated since. The eighteen months before it measure a mixture of judgement and memory of what
followed each headline.
\end{solution}

\begin{solution}{exo:ml:large-language-models-in-finance:2}
$\text{idf} = \log(1 + (540 - 45 + 0.5)/(45 + 0.5)) = \log 11.89 = 2.48$. At average length the normaliser is $k_1 = 1.5$, so
the term factor is $1\times2.5/(1 + 1.5) = 1$ and the contribution is 2.48.
\end{solution}

\begin{solution}{exo:ml:large-language-models-in-finance:3}
Premium (between minus after): contaminated $64.8 - 59.2 = 5.6$ points, clean $56.2 - 56.3 = -0.1$. Anonymisation loss
between the cut-offs: contaminated $64.8 - 57.0 = 7.8$, clean $56.2 - 56.0 = 0.2$.
\end{solution}

\begin{solution}{exo:ml:large-language-models-in-finance:4}
Its score depends only on the words it has embeddings for, the metric synonyms and context words; company names and years
are not in its vocabulary, so the 135 passages about the right metric tie and the top one is about the right company by
chance (1.1\%). Added to \omterm{def:ml:large-language-models-in-finance:rag}{BM25}, it lifts every passage about the right metric equally, including the other company's passage
that shares the question's synonym, which \omterm{def:ml:large-language-models-in-finance:rag}{BM25} already ranks first: the two scores disagree only where dense retrieval is
blind.
\end{solution}

\begin{solution}{exo:ml:large-language-models-in-finance:5}
It trained on a third more headlines (7\,988 to day 1\,210 instead of 5\,972 to day 907), so it estimates the events' effects better, and more
recently. After both cut-offs it has seen none of the answers, so the gain is generalisation, judgement in the chapter's
sense, bought with data. The point of the protocol is that this comparison is fair and the one between the cut-offs is not.
\end{solution}

\begin{solution}{exo:ml:large-language-models-in-finance:6}
Every headline is before the model's cut-off, so the 70\% mixes memory and judgement; the comparison with a tf-idf model
trained and tested out of sample is unfair. Removing the company names is not a full anonymisation (dates, products, people
and unique phrasing identify the event), and 65\% after it only shows that some memory survived or that the text is
informative, not which. Score on headlines after the cut-off, anonymise names and dates, and compare with baselines on the
same headlines.
\end{solution}

\begin{solution}{exo:ml:large-language-models-in-finance:7}
Changing the last token leaves the logits at every earlier position unchanged (the acceptance test checks it to $10^{-6}$).
The attention mask is upper-triangular with $-\infty$ above the diagonal, so after the softmax position $t$ puts zero weight
on positions after $t$; the feed-forward layers act position by position, so nothing else mixes positions.
\end{solution}

\begin{solution}{exo:ml:large-language-models-in-finance:8}
Question: the signals whose information coefficient fell most between the previous and the last quarter, with the
decay measure of Book~7 (chapter~13). Tools: the signal registry, the point-in-time data store with the as-of date set to the
quarter end, the IC calculator, and a plotting tool; no tool may write. Checks: the as-of guard, the minimum number of
observations per signal, a multiple-testing adjustment for ranking many signals. The log: task, model version, \omterm{def:ml:large-language-models-in-finance:prompt}{prompt},
every \omterm{def:ml:large-language-models-in-finance:agent}{tool call} with arguments and results' hashes, refusals, the final answer, and the code that recomputes it without the
agent.
\end{solution}

\begin{solution}{pb:ml:large-language-models-in-finance:1}
\textbf{Part I.}
\begin{enumerate}
\item Sentences of a week token, the company, the headline and the outcome after an arrow, for every headline before its
cut-off; asked by prompting up to the arrow and comparing the probabilities of ``up'' and ``down''.
\item 106\,681 parameters and 505 tokens. Whether a model memorises depends on how often it saw a document and on what keys
identify it, not only on size; the protocol (period split, anonymisation) measures memory in any model.
\item 61.6\%: the sign of the planted effect is right 66.0\% of the time on the 1\,462 headlines with an effect, and the
other 550 are coin flips. The reaction's noise (3\%) is large against the effects.
\item So that the placeholders are words the model has seen; otherwise the \omterm{def:ml:large-language-models-in-finance:memo}{anonymisation test} would measure the model's
reaction to unknown tokens, not the loss of its memory.
\end{enumerate}
\textbf{Part II.}
\begin{enumerate}[resume]
\item Clean 69.7\%, contaminated 65.4\%; anonymised 57.7\% and 57.2\%. Both memorised the outcomes of their training
headlines, keyed by week and company; without the keys they fall to their judgement.
\item Contaminated 64.8\%, clean 56.2\%: the contaminated model saw these headlines and their outcomes.
\item It takes the contaminated model to 57.0\% (a loss of 7.8 points) and the clean one to 56.0\% (0.2).
\item Lopez-Lira and Tang evaluated on headlines after GPT-4's cut-off; Glasserman and Lin found that anonymised headlines
did better in-sample, because general knowledge of the company distracted the model more than look-ahead helped it.
\end{enumerate}
\textbf{Part III.}
\begin{enumerate}[resume]
\item Same word: \omterm{def:ml:large-language-models-in-finance:rag}{BM25} 100\%, dense 0\%, hybrid 100\%, expanded \omterm{def:ml:large-language-models-in-finance:rag}{BM25} 86.3\%. Other word: 5.0\%, 1.1\%, 18.2\%, 89.1\%.
\item It turns away 11.9\% of the answerable questions, exactly those whose top passage named another company or year and
would have been answered wrongly; it abstains on all 300 unanswerable questions, which would otherwise all get a number.
\item It answers 80.8\% of the as-of questions from a filing published later, and is right 19.2\% of the time; with the
guard 100\%.
\item Each \omterm{def:ml:large-language-models-in-finance:agent}{tool call} with its arguments, the as-of date, the result or refusal, and the model version and \omterm{def:ml:large-language-models-in-finance:prompt}{prompt}.
\end{enumerate}
\textbf{Part IV.}
\begin{enumerate}[resume]
\item \emph{The model that knew.} \omterm{def:ml:large-language-models-in-finance:memo}{Memorisation} premium 5.6 points for the contaminated model and $-0.1$ for the clean one;
anonymisation loss 7.8 points and 0.2.
\item Ask for its cut-off and version history, results by period relative to the cut-off, an anonymised run, and a live or
post-cut-off trial against a dictionary baseline.
\item Extraction, tagging and summarisation with the source quoted, and code drafting with tests: outputs a person or a test
can check.
\item Material non-public information, client data and anything the firm would not publish.
\item It runs more experiments, faster, without memory of what it tried, and uses whatever data the tools return; without a
guard and a log it overfits and looks ahead at scale.
\item Pin the model version, \omterm{def:ml:large-language-models-in-finance:prompt}{prompt}, decoding settings and retrieval index; log every call; keep a code path that recomputes
the result without the model.
\item When the task is stable, labelled data exist, and prompting leaves accuracy or cost on the table; and when the model
can be hosted where the data may go.
\item Before its cut-off, memory and judgement together; only after it, judgement.
\end{enumerate}
\end{solution}

\begin{solution}{iq:ml:large-language-models-in-finance:1}
The date of the last text in the training data. Before it the model may have read the event and its aftermath, so a
backtest over that period has look-ahead bias that no data hygiene on the firm's side can remove.
\iqlookfor{look-ahead through the model's training data, and the remedy (post-cut-off evaluation).}
\end{solution}

\begin{solution}{iq:ml:large-language-models-in-finance:2}
Retrieve passages, put them in the \omterm{def:ml:large-language-models-in-finance:prompt}{prompt}, generate an answer grounded in them. Evaluate retrieval with labelled
question--passage pairs (recall at $k$), and generation given the right passage (faithfulness, exact match), so that
failures are attributed to the right stage.
\iqlookfor{the two stages and separate metrics for each.}
\end{solution}

\begin{solution}{iq:ml:large-language-models-in-finance:3}
Compare accuracy before and after the stated cut-off; anonymise names and dates and measure the loss; test on events the
model cannot know (post-cut-off, or synthetic ones); check whether it can complete the headline's text verbatim.
\iqlookfor{the period split and anonymisation, ideally a completion test.}
\end{solution}

\begin{solution}{iq:ml:large-language-models-in-finance:4}
\omterm{def:ml:large-language-models-in-finance:rag}{BM25} for names, tickers, numbers and exact terms; embeddings for paraphrase; in practice both, with the embedding used to
expand the query or rerank \omterm{def:ml:large-language-models-in-finance:rag}{BM25}'s candidates, measured on labelled questions from the desk's own documents.
\iqlookfor{what each misses, and a measured hybrid.}
\end{solution}

\begin{solution}{iq:ml:large-language-models-in-finance:5}
Put the as-of date in the harness, not the \omterm{def:ml:large-language-models-in-finance:prompt}{prompt}: every data tool filters by publication time and refuses later results,
and every call is logged. Prove it by the log and by a test that plants future documents and checks they are refused.
\iqlookfor{enforcement in the tools, a log, and a planted-future test.}
\end{solution}

\begin{solution}{iq:ml:large-language-models-in-finance:6}
The model and version behind it and their cut-offs; results split before and after the cut-off; an anonymised backtest; the
live history since launch; the timestamps of scores against news; costs and capacity; and a free trial on data after the
cut-off.
\iqlookfor{the cut-off split, anonymisation, live history and timestamps.}
\end{solution}
